\documentclass{article}
\usepackage[utf8]{inputenc}
\usepackage{amsmath,amssymb}
\usepackage[most]{tcolorbox}
\usepackage{geometry}
\geometry{margin=2.5cm}

\newcommand{\step}[1]{\par\noindent\hangindent=3.5em\hangafter=1%
  \makebox[3.5em][l]{\textbf{#1.}}}
\newcommand{\steptag}[1]{\hfill{\small\textsf{[#1]}}}

\newtcolorbox{proofbox}[1]{
  colback=gray!5, colframe=gray!60,
  fonttitle=\bfseries, title={#1},
  breakable, enhanced,
  left=4pt, right=4pt, top=4pt, bottom=4pt,
  before skip=10pt, after skip=10pt
}

\begin{document}

\begin{proofbox}{There are universal C\_pair<infinity and e\_pair>0 such tha...}
\step{1} There are universal C\_pair<infinity and e\_pair>0 such that every finite-dimensional extended epsilon\_X-C*-algebra (calX,I\_X,.,dagger) with 0<=epsilon\_X<=e\_pair and 1<dim\_C calX<infinity contains nonvanishing C\_pair*epsilon\_X-projections P',P'' with P'+P''=I\_X for which the linear map v\textasciicircum{}(2):C\textasciicircum{}2->calX, v\textasciicircum{}(2)(lambda,mu)=lambda*P'+mu*P'', is an extended C\_pair*epsilon\_X-inclusion, satisfies v\textasciicircum{}(2)(1,1)=I\_X, and sends the standard projection basis Pi',Pi'' to P',P''. \steptag{validated} \\[6pt]
\step{1.1} Fix an arbitrary finite-dimensional extended epsilon\_X-C*-algebra satisfying the root hypotheses. By lem-stage1-original-complementary-pair there are universal C\_np<infinity and e\_np>0 and, whenever epsilon\_X<=e\_np, elements q\_1,q\_2 with q\_1+q\_2=I\_X such that each q\_i is a nonvanishing C\_np*epsilon\_X-projection and ||q\_1q\_2||,||q\_2q\_1||<=C\_np*epsilon\_X. Thus, with d=max\{||q\_1\textasciicircum{}2-q\_1||,||q\_2\textasciicircum{}2-q\_2||,||q\_1q\_2||,||q\_2q\_1||\}, one has d<=C\_np*epsilon\_X and q\_i\textasciicircum{}dagger=q\_i. \steptag{validated} \\[4pt]
\step{1.2} By the nonvanishing alternative in def-delta-projection, applied to the pair from node 1.1, there is a universal e\_nv in (0,e\_np] such that epsilon\_X<=e\_nv implies ||q\_1||>=1/2 and ||q\_2||>=1/2. Indeed nonvanishing gives | ||q\_i||-1 |<=O(C\_np*epsilon\_X+epsilon\_X), with a universal data-independent coefficient, so one common shrink works for both i. \steptag{validated} \\[4pt]
\step{1.3} Define the single level-one linear map v\textasciicircum{}(2):C\textasciicircum{}2->calX by v\textasciicircum{}(2)(lambda,mu)=lambda*q\_1+mu*q\_2 and, for every n>=1, define only its canonical amplification v\_n=id\_\{M\_n\} tensor v\textasciicircum{}(2), so v\_n(A,B)=A tensor q\_1+B tensor q\_2. Linearity is immediate; q\_i\textasciicircum{}dagger=q\_i gives exact involution preservation; for the standard projection basis Pi'=(1,0), Pi''=(0,1) from def-projection-basis one has v\textasciicircum{}(2)(Pi')=q\_1 and v\textasciicircum{}(2)(Pi'')=q\_2; and q\_1+q\_2=I\_X gives v\textasciicircum{}(2)(1,1)=I\_X and v\_n(I\_n,I\_n)=I\_n tensor I\_X exactly. \steptag{validated} \\[4pt]
\step{1.4} The operator-space axioms in def-operator-space imply the simple-tensor identity ||T tensor z||\_n=||T||*||z|| for T in M\_n and z in calX: the rectangular matrix inequality gives the upper bound, while compression by unit vectors u,v chosen with |u\textasciicircum{}*Tv|=||T|| gives ||T||*||z||=||(u\textasciicircum{}* tensor 1)(T tensor z)(v tensor 1)||<=||T tensor z||. This applies uniformly at every n. \steptag{validated} \\[4pt]
\step{1.5} For x=(A,B), y=(C,D) in M\_n(C\textasciicircum{}2), bilinearity alone gives the exact four-term identity v\_n(x)v\_n(y)-v\_n(xy)=AC tensor(q\_1\textasciicircum{}2-q\_1)+AD tensor(q\_1q\_2)+BC tensor(q\_2q\_1)+BD tensor(q\_2\textasciicircum{}2-q\_2). Using node 1.4, ||RS||<=||R||*||S|| for scalar matrices, and ||(A,B)||=max\{||A||,||B||\}, its norm is at most 4d||x||||y||. Together with the exact linearity and involution preservation in node 1.3, every v\_n is a non-unital 4d-homomorphism, with a defect independent of n. \steptag{validated} \\[4pt]
\step{1.6} Assume epsilon\_X<=e\_nv and normalize ||(A,B)||=1. If ||A||=1, then v\_n(A,B)(I\_n tensor q\_1)-A tensor q\_1=A tensor(q\_1\textasciicircum{}2-q\_1)+B tensor(q\_2q\_1), so node 1.4 gives error at most 2d. Since M\_n tensor calX is an epsilon\_X-C*-algebra by def-extended-epsilon-cstar-algebra, its defining multiplication bound gives ||v\_n(A,B)(I\_n tensor q\_1)||<=(1+epsilon\_X)||v\_n(A,B)||*||q\_1||. Hence ||v\_n(A,B)||>=(||q\_1||-2d)/((1+epsilon\_X)||q\_1||)>=(1-4d)/(1+epsilon\_X). If instead ||B||=1, the same calculation with right multiplication by q\_2 and the defect q\_1q\_2 gives the identical bound. After a universal further shrink ensuring d<=1/8 and epsilon\_X<=1, homogeneity yields ||v\_n(x)||>=1/4||x|| for every n and x. \steptag{validated} \\[4pt]
\step{1.7} Choose a universal e\_up>0 no larger than all preceding thresholds, the smallness radii in GT-kitaev-prop-delta-hominc, and a number satisfying 8*C\_np*e\_up<1/4. Then delta\_n:=4d<=4*C\_np*epsilon\_X has 2*delta\_n<1/4. Apply GT-kitaev-prop-delta-hominc to each fixed canonical v\_n from node 1.3: the domain M\_n(C\textasciicircum{}2) is an exact C*-algebra (hence an epsilon\_X-C*-algebra), the codomain M\_n tensor calX is an epsilon\_X-C*-algebra by def-extended-epsilon-cstar-algebra, node 1.5 supplies the delta\_n-homomorphism property, and node 1.6 supplies eta=1/4. The external gives ||v\_n||<=1+O(delta\_n+epsilon\_X) and ||v\_n(x)||>=(1-O(delta\_n+epsilon\_X))||x||, with universal big-O coefficients independent of n and all input data. \steptag{validated} \\[6pt]
\step{1.7.1} Let delta\_max>0 be the universal delta-smallness radius in GT-kitaev-prop-delta-hominc, and put M:=max\{C\_np,1\}. In the choice already made at node 1.7, shrink the universal positive threshold e\_up further so that e\_up<=delta\_max/(4M) and e\_up<=1/(64M), as well as every preceding threshold. (This is possible because C\_np is finite and all listed radii are positive.) Since 0<=epsilon\_X<=e\_up and d<=C\_np*epsilon\_X, for every n we have delta\_n=4d<=4C\_np*epsilon\_X<=4M*e\_up<=delta\_max; independently, 2delta\_n<=8C\_np*e\_up<=8M*e\_up<=1/8<1/4. Thus both the separate hypothesis delta\_n<=delta\_max and the eta compatibility 2delta\_n<eta=1/4 required by GT-kitaev-prop-delta-hominc hold. Together with nodes 1.3, 1.5, and 1.6 and the domain/codomain facts stated in node 1.7, the proposition is therefore applicable to each v\_n and yields exactly the two uniform norm estimates asserted there. \steptag{validated} \\[4pt]
\step{1.8} Because delta\_n+epsilon\_X<=(4*C\_np+1)epsilon\_X, the universal estimates in node 1.7 admit one finite universal coefficient K such that, simultaneously for all n, v\_n is a K*epsilon\_X-homomorphism and (1-K*epsilon\_X)||x||<=||v\_n(x)||<=(1+K*epsilon\_X)||x||. By def-extended-delta-inclusion, the one map v\textasciicircum{}(2) is therefore an extended K*epsilon\_X-inclusion; no level-dependent map or constant was chosen. \steptag{validated} \\[6pt]
\step{1.8.1} Write the two universal big-O bounds supplied by node 1.7 as U*(delta\_n+epsilon\_X) for the upper estimate and L*(delta\_n+epsilon\_X) for the lower estimate, with fixed finite universal U,L independent of n and of the input algebra (enlarging either coefficient if necessary on the common smallness range). Since delta\_n=4d<=4*C\_np*epsilon\_X, choose K=max\{4*C\_np,U*(4*C\_np+1),L*(4*C\_np+1)\}. Then for every n and x one has the multiplicative defect at most 4d||x||||y||<=K*epsilon\_X||x||||y||, and (1-K*epsilon\_X)||x||<=||v\_n(x)||<=||v\_n||*||x||<=(1+K*epsilon\_X)||x||. This is one K for all n. This bridge is conditional on, and may be accepted only after, validation of node 1.7. \steptag{validated} \\[4pt]
\step{1.9} Set e\_pair=e\_up and choose the universal C\_pair>=max\{C\_np,K,4*C\_np\}. Nodes 1.1-1.2 make q\_1,q\_2 nonvanishing C\_pair*epsilon\_X-projections with q\_1+q\_2=I\_X (enlarging the defect allowance preserves that property); nodes 1.3 and 1.8 give the required single linear map, its extended C\_pair*epsilon\_X-inclusion bounds at every amplification, the exact unit value, and the exact standard-basis images. Taking P'=q\_1 and P''=q\_2 proves every clause of the root contract. \steptag{validated} \\[4pt]
\end{proofbox}

\end{document}
